\documentclass[11pt,letterpaper]{article}
\usepackage[T1]{fontenc}
\usepackage{lmodern}
\usepackage[margin=0.95in]{geometry}
\usepackage{amsmath,amssymb}
\usepackage{enumitem}
\usepackage[hidelinks]{hyperref}
\hypersetup{pdftitle={The SU(2) Wilson continuum limit and the Osterwalder-Schrader axioms}}
\setlength{\parindent}{0pt}
\setlength{\parskip}{5pt}
\setlist{leftmargin=*,itemsep=3pt,topsep=3pt,parsep=0pt}
\newcommand{\R}{\mathbb R}
\newcommand{\C}{\mathbb C}
\newcommand{\E}{\mathbb E}
\newcommand{\SU}{\mathrm{SU}}
\newcommand{\J}{\mathcal J}
\begin{document}
\begin{center}
{\LARGE A subsequential $\SU(2)$ Wilson continuum limit}\\[4pt]
{\large The Osterwalder--Schrader axioms for the limiting components}
\end{center}

\textbf{The observable.}
On a periodic four-dimensional lattice of mesh $a=L^{-K}$ and physical side
$\ell=2L^m$, give the $\SU(2)$ bond variables $U_b$ the Wilson probability law
\[
  d\mu_{K,\ell,g}=Z^{-1}\exp\!\left(-g^{-2}\sum_p s_p\right)\prod_b dU_b,
  \qquad s_p=1-\tfrac12\operatorname{Re}\operatorname{tr}U_p .
\]
Here $dU_b$ is normalized Haar measure, $p$ runs over positively oriented
plaquettes, and $U_p$ is the ordered product around $p$, taking the inverse on
a reversed bond. For a real compactly supported test $f$ on $\R^4$, define
\[
 f_\ell(x)=\sum_{z\in\mathbb Z^4}f(x+\ell z),\qquad
 \Phi(f)=\sum_p f_\ell(x_p)(s_p-\E s_p),
 \qquad
 M_n(K,\ell,g;\mathbf f)=\E\!\prod_{i=1}^n\Phi(f_i),
\]
where $x_p$ is the plaquette barycenter and expectations use $\mu_{K,\ell,g}$.

\textbf{The fixed measurement.}
Set $b(t)=e^{-1/t}$ for $t>0$ and $b(t)=0$ otherwise. With
$h(t)=b(t)/(b(t)+b(1-t))$, put
$\chi(x)=h(2-(80/3)^2|x|^2)$, and let $\chi_0(x)=\chi(x-3e_1)$,
where $e_1$ is the first coordinate vector. The bump $\chi$ equals one at the
origin and is supported in the ball of radius $1/8$. On a reference torus
of side $\ell_0$, set
\[
 F_K(g)=M_2(K,\ell_0,g;(\chi,\chi_0)),\qquad
 r(g_0)=\frac{9g_0^4}{\pi^4}\int_{\R^4}\!\int_{\R^4}
 \frac{\chi(x)\chi_0(y)}{|x-y|^8}\,dy\,dx .
\]

\textbf{Theorem.}
For every sufficiently large odd integer $L$, there exist integers $M\ge1$,
$K_*\ge1$ and a number $\delta>0$ with the following properties.
Take $\ell_0=2L^M\ge8$. For every $K\ge K_*$ and every $0<g_0<\delta$,
the equation $F_K(g)=r(g_0)$ has a least positive solution $g_K$.

Fix such a $g_0$. There are strictly increasing cutoffs $K_j$, observation
sides $\ell_j=2L^{m_j}\to\infty$, and, for each $j$, sides
$\ell_{j,t}=2L^{v(j,t)}\to\infty$ as $t\to\infty$. At each $K_j$ the
coupling $g_{K_j}$ is the least positive solution above. Asymptotic freedom takes the form
\[
 a_j=L^{-K_j}\longrightarrow0,\qquad
 g_{K_j}^{\,2}\log(1/a_j)\longrightarrow\frac{3\pi^2}{11}.
\]
There exists a family of Schwinger distributions $S_n$ on four-space
such that for every $n\ge0$ and all real $f_1,\ldots,f_n\in C_c^\infty(\R^4)$
with pairwise disjoint supports,
\[
 \lim_{j\to\infty}\lim_{t\to\infty} M_n(K_j,\ell_{j,t},g_{K_j};\mathbf f)
 =\lim_{j\to\infty}M_n(K_j,\ell_j,g_{K_j};\mathbf f)
 =S_n(f_1\otimes\cdots\otimes f_n).
\]
Each inner limit exists. The same sequences work at every order for every
separated test tuple. Some smooth separated pair has strictly positive
limiting covariance, and some smooth separated quadruple has strictly
positive limiting fourth cumulant.

The limiting moments are averages of components that satisfy the full
Osterwalder--Schrader axioms, including spatial clustering. Almost every
component has a nonzero connected four-point function on some separated
smooth quadruple. The test space, averaging formula and axioms are given below.

\newpage
\section*{Euclidean axioms and component decomposition}

Write $X_n=(\R^4)^n$ and $\Delta_n=\{x\in X_n:x_i=x_j\text{ for some }i\ne j\}$.
For $n\ge1$ the test space is
\[
 \mathcal S_n^{\flat}
 =\{\varphi\in\mathcal S(X_n;\C):
          \partial^\alpha\varphi|_{\Delta_n}=0\ \text{for every multi-index }\alpha\}.
\]
In degree zero take $\mathcal S_0^{\flat}=\C$. There are a probability space
$(\Xi,\rho)$ and component Schwinger distributions $T^\xi$ such that, for
every smooth separated tuple as on page~1, the integrand below is integrable and
\[
 S_n(f_1\otimes\cdots\otimes f_n)
       =\int_\Xi T^\xi_n(f_1\otimes\cdots\otimes f_n)\,d\rho(\xi).
\]
The following conditions hold for $S$ and for each $T^\xi$. Write $A$ for
either family.
\begin{enumerate}
\item \textbf{Normalization, linearity and growth.}
$A_0(c)=c$. Each $A_n$ is complex linear and continuous on $\mathcal S_n^{\flat}$.
For some integer $k>0$ and $c,C,\kappa>0$, every $n\ge1$ satisfies
\[
 |A_n(\varphi)|\le cC^n(n!)^\kappa p_{kn}(\varphi),\qquad
 p_d(\varphi)=\sup_{x,\,|\alpha|\le d}
             (1+|x|^2)^{d/2}|\partial^\alpha\varphi(x)|.
\]
\item \textbf{Euclidean invariance and permutation symmetry.}
Simultaneously translating all points, applying any $R\in O(4)$ to them,
or permuting the arguments leaves $A_n$ invariant.
\item \textbf{Reflection reality.} Write $\theta$ for reflection of the fourth
Euclidean coordinate and
$\J\varphi(x)=\overline{\varphi(\theta x_1,\ldots,\theta x_n)}$.
Then $A_n(\J\varphi)=\overline{A_n(\varphi)}$.
\item \textbf{Reflection positivity.} Let $F=(F_n)$ have only finitely many
nonzero entries. Take $F_0\in\C$, and, for $n\ge1$, take $F_n$ Schwartz with
all derivatives vanishing outside the set of distinct points whose time
coordinates are all positive. With $\J F_0=\overline{F_0}$, put
\[
 B_A(F,H)=\sum_{p,q\ge0} A_{p+q}(\J F_p\otimes H_q).
\]
Then $B_A(F,F)$ is real and nonnegative.
\end{enumerate}

\textbf{Clustering.} Each component $T^\xi$ also satisfies the spatial cluster
axiom. For $F,H$ as above and each nonzero spatial vector $v$,
\[
 \lim_{t\to+\infty} B_{T^\xi}(F,\tau_{tv} H)
   =\left(\sum_p T^\xi_p(\J F_p)\right)\left(\sum_q T^\xi_q(H_q)\right),
\]
where $\tau_w$ translates each test by $x_i\mapsto x_i-w$.
The separated smooth quadruple witnessing a nonzero fourth cumulant may
depend on $\xi$.

\textbf{The construction.} At each $a_j$ the inner limit on page~1 is a moment
of a probability law on fields on $a_j\mathbb Z^4$, obtained as a subsequential
weak limit of periodic Wilson measures. The construction also gives convergence
of separated connected correlations on compactly supported real $C^1$ tests,
and a canonical component decomposition on a compact metrizable space.

\textbf{Lean statement.} \texttt{YangMills.wilson\_subsequence} proves the
$\SU(2)$ instance. Its foundational axioms are \texttt{propext},
\texttt{Classical.choice}, and \texttt{Quot.sound}.
\end{document}
